\section*{Предисловие}
\addcontentsline{toc}{section}{Предисловие}
\label{sec:predislovie}

Курсовата работа е разработена по дисциплината „Електронни доказателства“ в Техническия
университет~-- София. Темата е формулирана от автора въз основа на въпроси 11, 12, 23 и 24
от конспекта на дисциплината: квалифицирания електронен подпис и механизма му на действие,
електронните доказателства в гражданския процес и представянето и оспорването им.

Изборът е свързан със специалността „Киберсигурност и превенция на киберпрестъпленията“.
Спорът дали определен човек е подписал документ с КЕП е правен, но се решава с технически
средства: валидиране на подписа, проверка на статута на удостоверението и анализ на
данните, които пази доставчикът. Работата е написана от студент с техническа подготовка и
затова съчетава правния и техническия поглед. Отделна част е посветена на казуса с липсата
на представителна власт, разгледан в курса.

Нормативните актове са цитирани по консолидираните им текстове към 16~ноември 2026~г.
Последното отразено изменение на ГПК е обнародвано в ДВ, бр.~92 от 2026~г., а ЗЕДЕУУ не е
изменян след ДВ, бр.~58 от 2019~г. Регламент (ЕС) №~910/2014 е цитиран по консолидирания
текст към 18.10.2024~г., който отразява измененията с Регламент (ЕС) 2024/1183.

Авторът благодари на доц. д-р Ралица Димитрова за материалите и казусите, предоставени в
курса.
